\section{Resolution behaviour of the slope}
\label{sec:slopesens}

Section~\ref{sec:validation} established that the framework reproduces
independent results when the mesh is fixed. This section asks a different
question: how much of a reported slope is a property of the body, and how much
is a property of the mesh on which it was computed.

\subsection{The measurement}
\label{sec:slopemeas}

For each body a single delivered shape model is decimated to a sequence of five
target facet counts spanning roughly $2\,000$ facets to the full model, using
the vertex-clustering operator of Section~\ref{sec:operators}. Density and
rotation period are held fixed at the values of Table~\ref{tab:shapes}, so the
body, its mass, and its spin are identical across the sequence; only the
sampling of its surface changes. The area-weighted mean slope $\slopeavg$ of
Eq.~\eqref{eq:slopeavg} is evaluated on each mesh.

Five bodies are used: Ryugu, Eros, Itokawa, Gaspra, and Ida. Kleopatra is
excluded, because both available models are already coarse ($4\,092$ and
$2\,292$ facets) and saturate under decimation, leaving too few distinct levels
to characterise a trend.

\subsection{Results}
\label{sec:slopesensres}

Table~\ref{tab:slopeseries} gives the full sequences, and the upper panel of
Figure~\ref{fig:resolution} plots them as fractional departures from the
finest-mesh value. On every body without exception $\slopeavg$ increases
monotonically with facet count. The increase is
not marginal: over the tested range it amounts to $3.5^{\circ}$ on Ryugu,
$2.2^{\circ}$ on Eros and $2.0^{\circ}$ on Itokawa, on bodies whose mean slopes
are themselves only $11$--$17^{\circ}$.

\begin{table}[t]
\centering
\small
\caption{Area-weighted mean slope $\slopeavg$ as a function of facet count for
a single shape model of each body, decimated by vertex clustering. Density and
rotation period are fixed at the values of Table~\ref{tab:shapes}.}
\label{tab:slopeseries}
\begin{tabular}{lrr}
\toprule
Body & $\Nf$ & $\slopeavg$ \\
\midrule
Ryugu   & $3\,096$  & $12.174^{\circ}$ \\
        & $5\,780$  & $12.829^{\circ}$ \\
        & $11\,134$ & $13.688^{\circ}$ \\
        & $20\,014$ & $14.505^{\circ}$ \\
        & $42\,834$ & $15.700^{\circ}$ \\
\midrule
Eros    & $2\,097$  & $11.034^{\circ}$ \\
        & $6\,306$  & $11.691^{\circ}$ \\
        & $12\,155$ & $12.122^{\circ}$ \\
        & $21\,850$ & $12.705^{\circ}$ \\
        & $36\,030$ & $13.263^{\circ}$ \\
\midrule
Itokawa & $2\,183$  & $14.710^{\circ}$ \\
        & $6\,404$  & $15.167^{\circ}$ \\
        & $12\,272$ & $15.762^{\circ}$ \\
        & $21\,196$ & $16.304^{\circ}$ \\
        & $41\,369$ & $16.727^{\circ}$ \\
\midrule
Gaspra  & $2\,821$  & $12.413^{\circ}$ \\
        & $5\,240$  & $12.466^{\circ}$ \\
        & $12\,500$ & $12.564^{\circ}$ \\
        & $20\,773$ & $12.619^{\circ}$ \\
        & $32\,040$ & $12.641^{\circ}$ \\
\midrule
Ida     & $2\,141$  & $13.171^{\circ}$ \\
        & $5\,808$  & $13.560^{\circ}$ \\
        & $11\,804$ & $13.807^{\circ}$ \\
        & $20\,673$ & $13.918^{\circ}$ \\
        & $32\,040$ & $13.933^{\circ}$ \\
\bottomrule
\end{tabular}
\end{table}

The two diagnostics of Section~\ref{sec:diagnostics} summarise these sequences
and, as anticipated there, they do not tell the same story
(Table~\ref{tab:slopesens}; see also the left-hand bars of each panel of
Figure~\ref{fig:diagnostics}).

The range sensitivity $S[\slopeavg]$ spans an order of magnitude between
bodies, from $1.8\%$ on Gaspra to $25.6\%$ on Ryugu. A quarter of the reported
value on Ryugu, and nearly a fifth on Eros, is attributable to the choice of
mesh alone.

The convergence residual $C[\slopeavg]$ is the more troubling of the two. On
Gaspra and Ida the slope has effectively settled by the finest mesh, changing
by only $0.17\%$ and $0.11\%$ over the last refinement step. On Ryugu, Eros and
Itokawa it has not: the last step still moves the value by $8.24\%$, $4.39\%$
and $2.59\%$ respectively, with no sign of a plateau. For those bodies there is
no evidence, within the resolution available, that $\slopeavg$ tends to a limit
at all. A value quoted at the finest mesh is therefore not an estimate of a
converged quantity; it is the value at that mesh.

\begin{table}[t]
\centering
\small
\caption{Range sensitivity $S$ and convergence residual $C$ of the
area-weighted mean slope, from the sequences of Table~\ref{tab:slopeseries}.
$C$ is the fractional change over the final refinement step.}
\label{tab:slopesens}
\begin{tabular}{lrrr}
\toprule
Body & $\Nf$ range & $S[\slopeavg]$ & $C[\slopeavg]$ \\
\midrule
Ryugu   & $3\,096$--$42\,834$ & $25.6\%$ & $8.24\%$ \\
Eros    & $2\,097$--$36\,030$ & $18.3\%$ & $4.39\%$ \\
Itokawa & $2\,183$--$41\,369$ & $12.8\%$ & $2.59\%$ \\
Ida     & $2\,141$--$32\,040$ & $5.6\%$  & $0.11\%$ \\
Gaspra  & $2\,821$--$32\,040$ & $1.8\%$  & $0.17\%$ \\
\bottomrule
\end{tabular}
\end{table}

\subsection{The origin of the trend}
\label{sec:slopeorigin}

That the slope rises with refinement, always in the same direction, points to a
specific mechanism. Decimation does not remove relief from the surface
(Section~\ref{sec:operators}); a coarse facet spans the same topography that
several fine facets spanned before, and its normal is close to the
area-weighted average of theirs. Averaging unit normals over a rough patch
yields a direction more nearly aligned with the mean surface than any individual
normal, so a coarse mesh systematically under-represents local departures from
the mean figure. Refining the mesh restores them, and the mean slope rises
accordingly.

If that account is right, then removing the relief at fixed facet count should
lower the slope by a comparable amount. The Laplacian smoothing operator of
Eq.~\eqref{eq:laplace} tests this directly: it suppresses short-wavelength
relief while leaving $\Nf$, the global figure, the mass and the spin unchanged.
Table~\ref{tab:smooth} gives the outcome for Itokawa and Eros.

\begin{table}[t]
\centering
\small
\caption{Effect of Laplacian smoothing at fixed facet count
($\Nf = 34\,234$ for Itokawa, $36\,030$ for Eros). The last column is the
fraction of facets steeper than $30^{\circ}$.}
\label{tab:smooth}
\begin{tabular}{lrrrr}
\toprule
Body & Iterations & $\slopeavg$ & Median & $>30^{\circ}$ \\
\midrule
Itokawa & $0$  & $15.81^{\circ}$ & $12.34^{\circ}$ & $6.5\%$ \\
        & $1$  & $14.82^{\circ}$ & $11.53^{\circ}$ & $4.9\%$ \\
        & $3$  & $14.00^{\circ}$ & $10.98^{\circ}$ & $4.3\%$ \\
        & $5$  & $13.56^{\circ}$ & $10.77^{\circ}$ & $3.9\%$ \\
        & $10$ & $12.90^{\circ}$ & $10.51^{\circ}$ & $3.4\%$ \\
\midrule
Eros    & $0$  & $13.65^{\circ}$ & $12.05^{\circ}$ & $3.5\%$ \\
        & $1$  & $12.23^{\circ}$ & $10.87^{\circ}$ & $1.6\%$ \\
        & $3$  & $11.35^{\circ}$ & $10.11^{\circ}$ & $0.8\%$ \\
        & $5$  & $10.92^{\circ}$ & $9.74^{\circ}$  & $0.6\%$ \\
        & $10$ & $10.30^{\circ}$ & $9.25^{\circ}$  & $0.2\%$ \\
\bottomrule
\end{tabular}
\end{table}

Smoothing lowers the mean slope by $2.9^{\circ}$ on Itokawa and $3.4^{\circ}$
on Eros, comparable to the $2.0^{\circ}$ and $2.2^{\circ}$ that refinement adds
over the tested range, and it acts most strongly on the steep tail: the fraction
of facets above $30^{\circ}$ falls by roughly half on Itokawa and by more than
an order of magnitude on Eros. The two operators move $\slopeavg$ in opposite
directions for the same underlying reason, and the mechanism is confirmed. What
the mesh resolution controls is how much of the surface relief the slope field
is able to see.

This also explains why the effect is a property of the shape model rather than
of the asteroid. Gaspra and Ida, whose models derive from limited flyby imaging
and carry little sub-facet relief to recover, show the smallest sensitivities
and have converged by the finest mesh. Ryugu, whose model is a dense
photogrammetric reconstruction rich in unresolved detail, shows the largest.
$S[\slopeavg]$ is therefore better read as a measure of how much relief a shape
model still contains below its nominal resolution than as a characteristic of
the body itself.

\subsection{Consequences for reporting}
\label{sec:slopereport}

Three consequences follow, and they are practical rather than conceptual.

First, a mean slope quoted without its facet count is not reproducible. On three
of the five bodies here the same shape model yields values differing by more
than $2^{\circ}$ depending on a choice that is rarely stated.

Second, mean slopes computed on different shape models of the same body are not
comparable, and neither are values from different studies unless both
resolutions are known and similar. Section~\ref{sec:crossversion} quantifies
this for a case where two independent models of one asteroid are available.

Third, and least comfortably, on the bodies with the richest shape models no
value within the tested range can be quoted as converged. Reporting $\slopeavg$ at the finest available
mesh is defensible only if the mesh is stated alongside it, and any comparison
must then be made at matched resolution. The next section shows that a closely
related quantity does not share this behaviour.
